% Part: first-order-logic
% Chapter: completeness
% Section: henkin-expansion

\documentclass[../../../include/open-logic-section]{subfiles}

\begin{document}

\olfileid[hi]{fol}{com}{hen}

\olsection{हेंकिन विस्तार}

\begin{explain}
पूर्णता प्रमेय को सिद्ध करने की एक कठिनाई यह है कि पूर्ण संगत
समुच्चय~$\Gamma$ से बनाई गई हमारी संरचना को~$\Gamma$ के सभी परिमाणित
!!{formula}s सत्य बनाने चाहिए। इसकी गारंटी के लिए हम लियोन हेंकिन की एक
युक्ति का उपयोग करते हैं। मूलतः इस युक्ति में भाषा को अनंततः अनेक
!!{constant}s से विस्तारित करना और प्रत्येक !!{formula} $!A(x)$, जिसमें एक
मुक्त !!{variable} हो, के लिए निम्न रूप का सूत्र जोड़ना है:
\iftag{prvEx}
      {$\lexists[x][!A(x)] \lif !A(c)$}
      {$\lnot\lforall[x][!A(x)] \lif \lnot !A(c)$},
जहाँ $c$ नए !!{constant}s में से एक है। जब हम~$\Gamma$ को संतुष्ट करने वाली
!!{structure} बनाएँगे, तो इससे यह गारंटी होगी कि
\iftag{prvEx}
{प्रत्येक सत्य अस्तित्ववाची वाक्य का कोई साक्षी}
{प्रत्येक असत्य सार्ववाची वाक्य का कोई प्रतिदृष्टांत}
नए अचरों में मौजूद हो।
\end{explain}

\begin{prop}
\ollabel{prop:lang-exp}
यदि $\Gamma$, $\Lang L$ में संगत है और $\Lang L'$ को~$\Lang L$ में
!!a{denumerable} समुच्चय के नए !!{constant}s $\Obj d_0$,
$\Obj d_1$, \dots जोड़कर प्राप्त किया गया है, तो $\Gamma$,~$\Lang L'$ में
संगत है।
\end{prop}

\begin{defn}[संतृप्त समुच्चय]
!!{formula}s का समुच्चय $\Gamma$, भाषा $\Lang {L}$ में,
\emph{संतृप्त} तभी और केवल तभी है जब प्रत्येक
!!{formula}~$!A(x) \in \Frm[L]$, जिसमें एक मुक्त !!{variable}~$x$ हो, के
लिए कोई !!a{constant}~$c \in \Lang{L}$ ऐसा हो कि
\iftag{prvEx}
      {$\lexists[x][!A(x)] \lif !A(c) \in \Gamma$}
      {$\lnot\lforall[x][!A(x)] \lif \lnot !A(c) \in \Gamma$}।
\end{defn}

अगली प्रमेय के प्रमाण में निम्न परिभाषा का उपयोग होगा।

\begin{defn}
\ollabel{defn:henkin-exp}
$\Lang L'$ को \olref{prop:lang-exp} जैसा मानें। एक गणन
$!A_0(x_0)$, $!A_1(x_1)$, \dots नियत करें, जिसमें वे सभी
!!{formula}s~$!A_i(x_i)$ हों जो~$\Lang L'$ के हैं और जिनमें एक चर ($x_i$)
मुक्त आता है। हम
!!{sentence}s~$!D_n$ को~$n$ पर आगमन द्वारा परिभाषित करते हैं।

$c_0$ को उन $\Obj d_i$ में पहला !!{constant} मानें जिन्हें हमने~$\Lang{L}$
में जोड़ा था और जो~$!A_0(x_0)$ में नहीं आता। यह मानते हुए कि $!D_0$,
\dots,~$!D_{n-1}$ पहले ही परिभाषित
हो चुके हैं, $c_n$ को नए !!{constant}s~$\Obj d_i$ में पहला ऐसा अचर मानें जो
न तो $!D_0$, \dots,~$!D_{n-1}$ में और न~$!A_n(x_n)$ में आता है।

अब $!D_{n}$ को निम्न !!{formula} मानें: \iftag{prvEx}
{$\lexists[x_{n}][!A_{n}(x_{n})] \lif
  !A_{n}(c_{n})$}{$\lnot\lforall[x_{n}][!A_{n}(x_{n})] \lif \lnot
  !A_{n}(c_{n})$}।
\end{defn}

\begin{lem}
\ollabel{lem:henkin}
प्रत्येक संगत समुच्चय~$\Gamma$ को किसी संतृप्त संगत समुच्चय~$\Gamma'$ तक
विस्तारित किया जा सकता है।
\end{lem}

\begin{proof}
वाक्यों का संगत समुच्चय~$\Gamma$, भाषा~$\Lang{L}$ में दिया हो, तो
!!a{denumerable} समुच्चय के नए !!{constant}s जोड़कर भाषा को~$\Lang{L'}$ में
विस्तारित करें। \olref{prop:lang-exp} से $\Gamma$ इस समृद्धतर भाषा में भी
संगत है। आगे, $!D_i$ को \olref{defn:henkin-exp} जैसा मानें। रखें
\begin{align*}
\Gamma_0 & = \Gamma \\
\Gamma_{n+1} & = \Gamma_n \cup \{!D_n \}
\end{align*}
अर्थात् $\Gamma_{n+1} = \Gamma \cup \{ !D_0, \dots, !D_n \}$, और
$\Gamma' = \bigcup_{n} \Gamma_n$ रखें। स्पष्टतः $\Gamma'$ संतृप्त है।

यदि $\Gamma'$ असंगत होता, तो किसी $n$ के लिए $\Gamma_n$ असंगत होता
(अभ्यास: समझाइए क्यों)। अतः $\Gamma'$ की संगतता दिखाने के लिए~$n$ पर आगमन
से यह दिखाना पर्याप्त है कि प्रत्येक समुच्चय~$\Gamma_n$ संगत है।

आगमन का आधार केवल यह कथन है कि $\Gamma_0 = \Gamma$ संगत है, जो प्रमेय की
परिकल्पना है। आगमन-चरण के लिए मान लें कि $\Gamma_{n}$ संगत है, पर
$\Gamma_{n+1} =
\Gamma_n \cup \{!D_n\}$ असंगत है। याद करें कि $!D_n$ यह है:
\iftag{prvEx}
{$\lexists[x_{n}][!A_{n}(x_n)] \lif !A_{n}(c_{n})$}
{$\lnot\lforall[x_{n}][!A_{n}(x_n)] \lif \lnot !A_{n}(c_{n})$},
जहाँ $!A_n(x_n)$, $\Lang{L'}$ का ऐसा !!a{formula} है जिसमें केवल
चर~$x_n$ मुक्त है। जिस प्रकार हमने~$c_n$ चुना है (देखें
\olref{defn:henkin-exp}), उससे $c_n$ न तो~$!A_n(x_n)$ में आता है, न
ही~$\Gamma_n$ में।

यदि $\Gamma_n \cup \{!D_n\}$ असंगत है, तो $\Gamma_n
\Proves \lnot !D_n$, और इसलिए निम्न दोनों सत्य हैं:
\iftag{prvEx}{
\[
\Gamma_n \Proves \lexists[x_n][!A_n(x_n)]
\qquad
\Gamma_n \Proves \lnot !A_n(c_n)
\]}{
\[
\Gamma_n \Proves \lnot\lforall[x_n][!A_n(x_n)]
\qquad
\Gamma_n \Proves !A_n(c_n)
\]}
चूँकि $c_n$, $\Gamma_n$ या~$!A_n(x_n)$ में नहीं आता,
\tagrefs{prfAX/{fol:axd:qpr:thm:strong-generalization},prfSC/{fol:seq:qpr:thm:strong-generalization},prfND/{fol:ntd:qpr:thm:strong-generalization},prfTab/{fol:tab:qpr:thm:strong-generalization}} लागू होता है।
\iftag{prvEx}
{$\Gamma_n \Proves \lnot !A_n(c_n)$}
{$\Gamma_n \Proves !A_n(c_n)$}
से हमें
\iftag{prvEx}
{$\Gamma_n \Proves \lforall[x_n][\lnot !A_n(x_n)]$}
{$\Gamma_n \Proves \lforall[x_n][!A_n(x_n)]$}
मिलता है। अतः हमारे पास निम्न दोनों हैं:
\iftag{prvEx}
{$\Gamma_n \Proves \lexists[x_n][!A_n(x_n)]$ और
$\Gamma_n \Proves \lforall[x_n][\lnot !A_n(x_n)]$}
{$\Gamma_n \Proves \lnot\lforall[x_n][!A_n(x_n)]$ और
$\Gamma_n \Proves \lforall[x_n][!A_n(x_n)]$},
इसलिए $\Gamma_n$ स्वयं असंगत है।
\iftag{prvEx}
{(ध्यान दें कि
\iftag{prvAll}
{$\lforall[x_n][\lnot !A_n(x_n)] \Proves
  \lnot\lexists[x_n][!A_n(x_n)]$}
{$\lforall[x_n][\lnot !A_n]$ को
  $\lnot\lexists[x_n][\lnot\lnot !A_n(x_n)]$ के रूप में परिभाषित किया गया
  है और $\lnot\lexists[x_n][\lnot\lnot !A_n(x_n)] \Proves
  \lnot\lexists[x_n][!A_n(x_n)]$}।)}{}
विरोधाभास: $\Gamma_n$ को संगत माना गया था। अतः
$\Gamma_n \cup \{ !D_n\}$ संगत है।
\end{proof}

\begin{explain}
अब हम दिखाएँगे कि संतृप्त \emph{पूर्ण} संगत समुच्चयों में यह गुण होता है कि
\iftag{prvAll}{वे किसी सार्ववाची !!{sentence} को तभी और केवल तभी समाहित
  करते हैं जब वे उसके सभी निदर्शों को समाहित करें}{}\iftag{defAll,defEx}{}{
  और }\iftag{prvAll}{वे किसी अस्तित्ववाची !!{sentence} को तभी और केवल तभी
  समाहित करते हैं जब वे कम-से-कम एक निदर्श को समाहित करें}{}। इसका उपयोग हम
यह दिखाने के लिए करेंगे कि पूर्ण, संगत, संतृप्त समुच्चय से बनाई गई
!!{structure} उसके सभी परिमाणित वाक्यों को सत्य बनाती है।
\end{explain}

\begin{prop}\ollabel{prop:saturated-instances}
मान लें कि $\Gamma$ पूर्ण, संगत और संतृप्त है।
\begin{tagenumerate}{prvEx,prvAll}
\tagitem{prvEx}{$\lexists[x][!A(x)] \in \Gamma$ तभी और केवल तभी जब
  $!A(t) \in \Gamma$ कम-से-कम एक बंद पद~$t$ के लिए हो।}{}
\tagitem{prvAll}{$\lforall[x][!A(x)] \in \Gamma$ तभी और केवल तभी जब
  $!A(t) \in \Gamma$ सभी बंद पदों~$t$ के लिए हो।}{}
\end{tagenumerate}
\end{prop}

\begin{proof}
  \begin{tagenumerate}{prvEx,prvAll}
  \tagitem{prvEx}{%
    \iftag{probEx}{अभ्यास।}{पहले मान लें कि $\lexists[x][!A(x)]
      \in \Gamma$। क्योंकि $\Gamma$ संतृप्त है,
      $(\lexists[x][!A(x)] \lif !A(c)) \in \Gamma$ किसी
      !!{constant}~$c$ के लिए।
      \tagrefs{prfAX/{fol:axd:ppr:prop:provability-lif},prfSC/{fol:seq:ppr:prop:provability-lif},prfND/{fol:ntd:ppr:prop:provability-lif},prfTab/{fol:tab:ppr:prop:provability-lif}} के पद~(1)
      और \olref[ccs]{prop:ccs}\olref[ccs]{prop:ccs-prov-in} से $!A(c)
      \in \Gamma$।

    दूसरी दिशा में संतृप्ति आवश्यक नहीं है: मान लें कि
    $!A(t) \in \Gamma$। तब
    \tagrefs{prfAX/{fol:axd:qpr:prop:provability-quantifiers},prfSC/{fol:seq:qpr:prop:provability-quantifiers},prfND/{fol:ntd:qpr:prop:provability-quantifiers},prfTab/{fol:tab:qpr:prop:provability-quantifiers}} के पद~(1) से $\Gamma \Proves \lexists[x][!A(x)]$।
    \olref[ccs]{prop:ccs}\olref[ccs]{prop:ccs-prov-in} से
    $\lexists[x][!A(x)] \in \Gamma$।}}{}
  \tagitem{prvAll}{%
    \iftag{probAll}{अभ्यास।}{मान लें कि $!A(t) \in \Gamma$ सभी बंद
      पदों~$t$ के लिए। विरोधाभास के लिए मान लें कि
      $\lforall[x][!A(x)] \notin \Gamma$। चूँकि $\Gamma$ पूर्ण है,
      $\lnot\lforall[x][!A(x)] \in \Gamma$। संतृप्ति से,
      \iftag{prvEx}{$(\lexists[x][\lnot !A(x)] \lif \lnot !A(c)) \in
        \Gamma$}{$(\lnot\lforall[x][!A(x)] \lif \lnot !A(c)) \in
        \Gamma$} किसी !!{constant}~$c$ के लिए। परिकल्पना से, क्योंकि $c$ बंद पद है,
      $!A(c) \in \Gamma$। किंतु इससे $\Gamma$ असंगत हो जाता। (अभ्यास: वह
      !!{derivation} दीजिए जो दिखाती है कि
      \[
      \lnot \lforall[x][!A(x)],
      \iftag{prvEx}{\lexists[x][\lnot !A(x)] \lif \lnot
        !A(c)}{\lnot\lforall[x][!A(x)] \lif \lnot
        !A(c)}, !A(c)
      \]
      असंगत है।)

      विपरीत दिशा के लिए संतृप्ति की आवश्यकता नहीं है: मान लें कि
      $\lforall[x][!A(x)] \in \Gamma$। तब
      \tagrefs{prfAX/{fol:axd:qpr:prop:provability-quantifiers},prfSC/{fol:seq:qpr:prop:provability-quantifiers},prfND/{fol:ntd:qpr:prop:provability-quantifiers},prfTab/{fol:tab:qpr:prop:provability-quantifiers}}
      के पद~(2) से $\Gamma \Proves !A(t)$। \olref[ccs]{prop:ccs} से हमें
      $!A(t) \in \Gamma$ मिलता है।}}{}
  \end{tagenumerate}
\end{proof}

\end{document}
